\begin{table}[t]\centering\caption{Main results: final validation loss (nats, mean $\pm$ std over \rhval{const/n_seeds} seeds) at the three main peak learning rates ($10^{-3}$, $3\times10^{-3}$, $10^{-2}$) and the two extra rates ($3\times10^{-4}$, $3\times10^{-2}$). Constant and step decay have no warmup here.}\label{tab:main}
\resizebox{\columnwidth}{!}{\begin{tabular}{llcccc}
\toprule
Method & Task & val\_loss $\downarrow$ & train\_loss $\downarrow$ & val\_loss\_y $\downarrow$ & $n$ \\
\midrule
Step decay & lr3e-3 & 1.85 $\pm$ 0.05 & 1.70 $\pm$ 0.05 & -- & 3 \\
Schedule-free AdamW & lr3e-3 & \underline{1.69 $\pm$ 0.01} & \underline{1.53 $\pm$ 0.00} & \textbf{1.72 $\pm$ 0.01} & 3 \\
Constant & lr3e-3 & 1.82 $\pm$ 0.03 & 1.66 $\pm$ 0.03 & -- & 3 \\
\textbf{Warmup+Cosine} & lr3e-3 & \textbf{1.67 $\pm$ 0.01} & \textbf{1.49 $\pm$ 0.00} & -- & 3 \\
\midrule
Step decay & lr1e-3 & 1.83 $\pm$ 0.00 & 1.68 $\pm$ 0.01 & -- & 3 \\
Schedule-free AdamW & lr1e-3 & 1.88 $\pm$ 0.01 & 1.75 $\pm$ 0.01 & \textbf{1.88 $\pm$ 0.01} & 3 \\
Constant & lr1e-3 & \textbf{1.76 $\pm$ 0.01} & \textbf{1.60 $\pm$ 0.00} & -- & 3 \\
\textbf{Warmup+Cosine} & lr1e-3 & \underline{1.77 $\pm$ 0.01} & \underline{1.61 $\pm$ 0.01} & -- & 3 \\
\midrule
Step decay & lr1e-2 & 2.35 $\pm$ 0.06 & 2.32 $\pm$ 0.05 & -- & 3 \\
Schedule-free AdamW & lr1e-2 & \textbf{1.63 $\pm$ 0.00} & \textbf{1.49 $\pm$ 0.00} & \textbf{1.67 $\pm$ 0.01} & 3 \\
Constant & lr1e-2 & 2.42 $\pm$ 0.05 & 2.40 $\pm$ 0.04 & -- & 3 \\
\textbf{Warmup+Cosine} & lr1e-2 & \underline{1.68 $\pm$ 0.01} & \underline{1.50 $\pm$ 0.00} & -- & 3 \\
\midrule
Step decay & lr3e-4 & 2.13 $\pm$ 0.01 & 2.09 $\pm$ 0.01 & -- & 3 \\
Schedule-free AdamW & lr3e-4 & 2.17 $\pm$ 0.01 & 2.13 $\pm$ 0.01 & \textbf{2.17 $\pm$ 0.01} & 3 \\
Constant & lr3e-4 & \textbf{1.94 $\pm$ 0.02} & \textbf{1.83 $\pm$ 0.01} & -- & 3 \\
\textbf{Warmup+Cosine} & lr3e-4 & \underline{2.07 $\pm$ 0.01} & \underline{2.02 $\pm$ 0.01} & -- & 3 \\
\midrule
Step decay & lr3e-2 & 2.50 $\pm$ 0.00 & 2.48 $\pm$ 0.00 & -- & 3 \\
Schedule-free AdamW & lr3e-2 & \textbf{1.65 $\pm$ 0.02} & \underline{1.55 $\pm$ 0.00} & \textbf{1.71 $\pm$ 0.01} & 3 \\
Constant & lr3e-2 & 2.80 $\pm$ 0.05 & 2.76 $\pm$ 0.03 & -- & 3 \\
\textbf{Warmup+Cosine} & lr3e-2 & \underline{1.71 $\pm$ 0.01} & \textbf{1.54 $\pm$ 0.00} & -- & 3 \\
\midrule
\bottomrule
\end{tabular}
}\end{table}

\begin{figure}[t]\centering\includegraphics[width=\columnwidth]{val_vs_lr.pdf}
\caption{Final validation loss against peak learning rate (log axis), mean and std over \rhval{const/n_seeds} seeds, five rates (group main).}\label{fig:lr}\end{figure}

\paragraph{H1 (cosine beats constant at the best LR): supported, but conditional.} Taking each system's best of the three main rates, warmup+cosine reaches \rhval{main/warmup+cosine/lr3e-3/val_loss/mean:2} at $3\times10^{-3}$, constant LR \rhval{main/constant/lr1e-3/val_loss/mean:2} at $10^{-3}$, step decay \rhval{main/step-decay/lr1e-3/val_loss/mean:2} at $10^{-3}$ and Schedule-Free \rhval{main/schedule-free-adamw/lr1e-2/val_loss/mean:2} at $10^{-2}$ (Table~\ref{tab:main}). The comparison of per-seed best losses, warmup+cosine minus constant, is \rhval{cmp/main/constant/sens/best3/delta:3} nats (Welch $p=\rhval{cmp/main/constant/sens/best3/welch_p}$, $n=\rhval{const/n_seeds}$ per arm; this follows the test registered in \texttt{research.yaml}). The advantage is not uniform across peak rates: at $10^{-3}$ the two are indistinguishable ($\Delta=\rhval{cmp/main/constant/lr1e-3/val_loss/delta:3}$, $p=\rhval{cmp/main/constant/lr1e-3/val_loss/welch_p}$), and at the lowest rate $3\times10^{-4}$ constant LR is \emph{better} than warmup+cosine (\rhval{main/constant/lr3e-4/val_loss/mean:2} vs.\ \rhval{main/warmup+cosine/lr3e-4/val_loss/mean:2}; $p=\rhval{cmp/main/constant/lr3e-4/val_loss/welch_p}$): when the peak rate is too small for a 1,500-step budget, decaying it only removes useful progress. At $3\times10^{-3}$ and $10^{-2}$ cosine wins by \rhval{cmp/main/constant/lr3e-3/val_loss/delta:2} and \rhval{cmp/main/constant/lr1e-2/val_loss/delta:2} nats respectively, but the $10^{-2}$ gap is largely a warmup effect (below).

\paragraph{Schedule-Free beats cosine at its best rate, but not everywhere.} At $10^{-2}$ Schedule-Free AdamW is lower than warmup+cosine (cosine minus Schedule-Free: \rhval{cmp/main/schedule-free-adamw/lr1e-2/val_loss/delta:3} nats; Welch $p=\rhval{cmp/main/schedule-free-adamw/lr1e-2/val_loss/welch_p}$) and likewise at $3\times10^{-2}$ (\rhval{main/schedule-free-adamw/lr3e-2/val_loss/mean:2} vs.\ \rhval{main/warmup+cosine/lr3e-2/val_loss/mean:2}). At $3\times10^{-3}$ the two are close (\rhval{main/schedule-free-adamw/lr3e-3/val_loss/mean:2} vs.\ \rhval{main/warmup+cosine/lr3e-3/val_loss/mean:2}, Welch $p=\rhval{cmp/main/schedule-free-adamw/lr3e-3/val_loss/welch_p}$), and at lower rates Schedule-Free is clearly worse: \rhval{main/schedule-free-adamw/lr1e-3/val_loss/mean:2} at $10^{-3}$ and \rhval{main/schedule-free-adamw/lr3e-4/val_loss/mean:2} at $3\times10^{-4}$, below even the constant schedule at those rates. Its loss at the non-averaged iterate $y_T$ is worse than at the average at the best rates (e.g.\ \rhval{main/schedule-free-adamw/lr1e-2/val_loss_y/mean:2} at $10^{-2}$), consistent with the averaging doing useful work.

\begin{table}[t]\centering\caption{Sensitivity: worst-minus-best validation loss over three ($\text{sens}_3$) and five ($\text{sens}_5$) peak rates, and best three-rate loss (best$_3$), mean $\pm$ std over seeds (lower is better).}\label{tab:sens}
\resizebox{\columnwidth}{!}{\begin{tabular}{llcccc}
\toprule
Method & Task & sens3 $\downarrow$ & sens5 $\downarrow$ & best3 $\downarrow$ & $n$ \\
\midrule
\midrule
\midrule
\midrule
\midrule
\midrule
Step decay & sens & 0.53 $\pm$ 0.05 & 0.68 $\pm$ 0.02 & 1.82 $\pm$ 0.02 & 3 \\
Schedule-free AdamW & sens & \underline{0.25 $\pm$ 0.01} & \underline{0.54 $\pm$ 0.01} & \textbf{1.63 $\pm$ 0.00} & 3 \\
Constant & sens & 0.65 $\pm$ 0.05 & 1.03 $\pm$ 0.03 & 1.76 $\pm$ 0.01 & 3 \\
\textbf{Warmup+Cosine} & sens & \textbf{0.09 $\pm$ 0.01} & \textbf{0.40 $\pm$ 0.02} & \underline{1.67 $\pm$ 0.01} & 3 \\
\bottomrule
\end{tabular}
}\end{table}

\paragraph{H2 (Schedule-Free less LR-sensitive than cosine): refuted.} Table~\ref{tab:sens} gives $\text{sens}_3$ of \rhval{main/schedule-free-adamw/sens/sens3/mean:2} for Schedule-Free and \rhval{main/warmup+cosine/sens/sens3/mean:2} for warmup+cosine (difference cosine minus Schedule-Free \rhval{cmp/main/schedule-free-adamw/sens/sens3/delta:2}, Welch $p=\rhval{cmp/main/schedule-free-adamw/sens/sens3/welch_p}$); over five rates the values are \rhval{main/schedule-free-adamw/sens/sens5/mean:2} and \rhval{main/warmup+cosine/sens/sens5/mean:2} ($p=\rhval{cmp/main/schedule-free-adamw/sens/sens5/welch_p}$). The direction is opposite to H2. Figure~\ref{fig:lr} shows why: Schedule-Free is flat and good for $10^{-2}$ and $3\times10^{-2}$ but degrades quickly below $3\times10^{-3}$, whereas cosine has a broad, shallow basin around $3\times10^{-3}$. Warmup+cosine is the least sensitive of the four \emph{as specified} (\rhval{main/constant/sens/sens3/mean:2} for constant, \rhval{main/step-decay/sens/sens3/mean:2} for step decay over three rates), but this ordering changes once the baselines get warmup (next section). Note that sensitivity is a range on a coarse grid: Schedule-Free's best loss is at or beyond the upper end of the main grid, so a grid shifted upward would make its $\text{sens}_3$ smaller; our result is about this grid, not about an intrinsic robustness property.

\begin{figure}[t]\centering\includegraphics[width=\columnwidth]{curves_lr1e-2.pdf}
\caption{Validation loss during training at peak LR $10^{-2}$ (seed 3, one run per system; evaluated every \rhval{const/eval_every} steps; Schedule-Free evaluated at the averaged iterate).}\label{fig:curves}\end{figure}

\paragraph{Training dynamics at $10^{-2}$.} Figure~\ref{fig:curves} shows that without warmup, constant and step decay (identical until the first drop) stay far above the schedules with warmup throughout training (final validation loss \rhval{curves/constant/lr1e-2/val_loss/mean:2} for constant and \rhval{curves/step-decay/lr1e-2/val_loss/mean:2} for step decay, against \rhval{curves/warmup+cosine/lr1e-2/val_loss/mean:2} for cosine and \rhval{curves/schedule-free-adamw/lr1e-2/val_loss/mean:2} for Schedule-Free in these runs). This is single-seed evidence for the early-training failure; it motivates, but does not by itself establish, the warmup explanation tested in the next section.

\paragraph{Failure cases and per-regime summary.} (i) Constant and step decay at $10^{-2}$ and above: \rhval{main/constant/lr1e-2/val_loss/mean:2} and \rhval{main/step-decay/lr1e-2/val_loss/mean:2} nats, \rhval{main/constant/lr3e-2/val_loss/mean:2} and \rhval{main/step-decay/lr3e-2/val_loss/mean:2} at $3\times10^{-2}$; no run was flagged as diverged, they simply train badly. (ii) Warmup+cosine at $3\times10^{-4}$ (\rhval{main/warmup+cosine/lr3e-4/val_loss/mean:2}), where it is worse than constant. (iii) Schedule-Free at $\le10^{-3}$. No system is best at all rates; the ranking depends on the regime, which is the main message of the sweep table.
